\section*{ЗАКЛЮЧЕНИЕ}
\addcontentsline{toc}{section}{ЗАКЛЮЧЕНИЕ}

Разработанная программно-информационная система обнаружения дефектов дорожного покрытия на основе методов компьютерного зрения и нейросетевой обработки изображений обеспечивает автоматический анализ изображений дорожного покрытия, обнаружение дефектов с использованием обученной модели YOLOv8n, визуализацию результатов обработки, сохранение обработанных изображений и формирование текстового отчёта.

Основные результаты работы:

\begin{enumerate}
	\item Проведён анализ предметной области и существующих методов обнаружения дефектов дорожного покрытия.
	\item Изучены особенности архитектуры нейросетевых моделей семейства YOLO.
	\item Выполнено проектирование программно-информационной системы и её пользовательского интерфейса.
	\item Реализована программная система обнаружения дефектов дорожного покрытия на языке Python.
	\item Выполнено обучение нейросетевой модели на специализированном наборе данных.
	\item Реализованы механизмы обработки изображений, сохранения результатов и формирования отчётов.
\end{enumerate}

Все требования, сформулированные в техническом задании, были реализованы в полном объёме. Поставленная цель разработки достигнута, все задачи, определённые в начале работы, успешно решены.

Результатом работы является программно-информационная система, позволяющая автоматизировать процесс обнаружения дефектов дорожного покрытия и сократить затраты времени на анализ изображений и сбор статистики по сравнению с ручными методами контроля.